% Part: first-order-logic
% Chapter: models-theories
% Section: theories

\documentclass[../../../include/open-logic-section]{subfiles}

\begin{document}

\olfileid{fol}{mat}{the}

\olsection{Tuladha Teori Logika Tataran Kapisan}

\begin{ex}
Teori urutan linear ketat ing basa~$\Lang L_<$ diaksiomatisasi
dening himpunan
\begin{align*}
\{\quad & \lforall[x][\lnot x < x], \\
& \lforall[x][\lforall[y][((x < y \lor y <
    x) \lor x = y)]], \\
& \lforall[x][\lforall[y][\lforall[z][((x < y
      \land y < z) \lif x < z)]]] \quad \}
\end{align*}
Iki nyakup kanthi pas !!{structure}s kang dikarepake:
saben urutan linear ketat ing himpunan kang ora kosong iku model saka sistem aksioma iki,
lan kosok balene, yen $R$ iku urutan linear ketat ing himpunan $X$ kang ora kosong,
mula struktur $\Struct M$ kanthi $\Domain M = X$ lan
$\Assign{<}{M} = R$ iku model saka teori iki. Cathetan koreksi sumber OLPL-752: relasi R ing arah kosok balen kudu urutan linear ketat supaya aksioma irrefleksif kasukupan. Cathetan cakupan SOL6-F008: himpunan dhasar uga kudu ora kosong amarga definisi struktur ora ngidini domain kosong.
\end{ex}

\begin{ex}
Teori grup ing basa kanthi $\Obj 1$ (!!{constant}) lan $\cdot$
(!!{function} loro panggonan) diaksiomatisasi dening
\begin{align*}
& \lforall[x][\eq[(x \cdot \Obj 1)][x]]\\
& \lforall[x][\lforall[y][\lforall[z][\eq[(x \cdot (y \cdot z))][((x
          \cdot y) \cdot z)]]]]\\
& \lforall[x][\lexists[y][\eq[(x \cdot y)][\Obj 1]]]
\end{align*}
\end{ex}

\begin{ex}
Teori aritmetika Peano diaksiomatisasi dening
ukara-ukara ing ngisor iki ing basa aritmetika~$\Lang L_A$.
\begin{align*}
& \lforall[x][\lforall[y][(\eq[x'][y'] \lif \eq[x][y])]]\\
& \lforall[x][\eq/[\Obj 0][x']]\\
& \lforall[x][\eq[(x + \Obj 0)][x]]\\
& \lforall[x][\lforall[y][\eq[(x + y')][(x + y)']]]\\
& \lforall[x][\eq[(x \times \Obj 0)][\Obj 0]]\\
& \lforall[x][\lforall[y][\eq[(x \times y')][((x \times y) + x)]]]\\
& \lforall[x][\lforall[y][(x < y \liff \lexists[z][\eq[(z' + x)][y])]]]\\
\intertext{ditambah kabeh ukara kang awujud}
& (!A(\Obj 0) \land \lforall[x][(!A(x) \lif !A(x'))]) \lif \lforall[x][!A(x)]
\end{align*}
Amarga ukara kanthi wujud kang pungkasan iku tanpa wates akehe,
sistem aksioma iki tanpa wates. Wujud kang pungkasan kasebut
\emph{skema induksi}. (Satemene, skema induksi rada luwih
rumit tinimbang kang wis kita andharake ing kene.)

Aksioma kang dumunung sadurunge skema induksi iku
\emph{definisi eksplisit} saka~$<$.
Cathetan koreksi sumber OLPL-083: ukara sumber nyebut
"aksioma pungkasan" sanajan skema induksi wis ditampilake
sawisé aksioma kang netepake relasi iki.
\end{ex}

\begin{ex}
Teori himpunan murni nduweni peran wigati ing dhasar
(lan ing filsafat) matematika. Himpunan diarani murni
yen kabeh !!{element}s uga himpunan murni. Mula,
himpunan kosong kalebu murni, nanging himpunan kang
nduweni sawenehing barang minangka !!a{element}
kang dudu himpunan ora murni. Dadi, himpunan murni
yaiku himpunan kang kawangun mung saka himpunan kosong
lan tanpa "urelemen", yaiku objek kang dudu himpunan.

Ukara-ukara ing ngisor iki bisa dianggep sistem aksioma
kanggo teori himpunan murni:
\begin{align*}
& \lexists[x][\lnot \lexists[y][y \in x]]\\
& \lforall[x][\lforall[y][(\lforall[z](z \in x \liff z \in y) \lif 
\eq[x][y])]]\\
& \lforall[x][\lforall[y][\lexists[z][\lforall[u][(u \in z \liff
(\eq[u][x] \lor \eq[u][y]))]]]]\\
& \lforall[x][\lexists[y][\lforall[z][(z \in y \liff \lexists[u][(z \in
        u \land u \in x)])]]]\\
\intertext{ditambah kabeh ukara kang awujud} &
\lexists[x][\lforall[y][(y \in x \liff !A(y))]]
\end{align*}
Aksioma kapisan nyatakake yen ana himpunan tanpa
!!{element}s (yaiku, $\emptyset$ ana); aksioma kapindho
nyatakake yen himpunan iku ekstensional; aksioma katelu
nyatakake yen kanggo sembarang himpunan $X$ lan $Y$,
himpunan $\{X, Y\}$ ana; aksioma kapapat nyatakake
yen kanggo sembarang himpunan $X$, himpunan $\cup X$ ana,
ing ngendi $\cup X$ iku gabungan kabeh anggota $X$.

!!{sentence}s kang kasebut pungkasan iku sacara bebarengan
diarani \emph{skema komprehensi naif}. Ing pokoke,
skema iki nyatakake yen kanggo saben $!A(x)$, himpunan
$\Setabs{x}{!A(x)}$ ana---mula nalika sepisanan dideleng,
iki katon minangka aksioma kang bener, migunani, lan
mbokmenawa malah perlu. Diarani "naif" amarga, kaya
kang bakal kita deleng, skema iki ndadekake teori
ora bisa dicukupi: yen $!A(y)$ dijupuk minangka
$\lnot y \in y$, kita oleh !!{sentence}
\[
\lexists[x][\lforall[y][(y \in x \liff \lnot y \in y)]]
\]
lan !!{sentence} iki ora dicukupi dening !!{structure} apa wae.
\end{ex}

\begin{ex}
Ing bidang \emph{mereologi}, relasi \emph{dadi perangan saka}
iku relasi dhasar. Kaya dene teori himpunan, ana teori
perangan kang ngaksiomatisasi maneka pamahaman
(kang kadhang padha nentang) ngenani relasi iki.

Basa mereologi ngemot siji simbol predikat loro panggonan,
yaiku~$\Obj P$, lan $\Atom{\Obj P}{x, y}$ "tegese" $x$
iku perangan saka~$y$. Nalika interpretasi iki kang kita
karepake, !!a{structure} kanggo basa iki diarani
\emph{struktur perangan}. Mesthi ora saben struktur
kanggo siji predikat loro panggonan pancen pantes
diarani mangkono. Supaya bisa nyakup makna
"dadi perangan saka", $\Assign{\Obj P}{M}$ kudu
nyukupi sawenehing syarat, kang bisa kita tetepake
minangka aksioma kanggo teori perangan. Upamane,
relasi perangan iku urutan parsial tumrap objek:
saben objek iku perangan (sanajan perangan
\emph{ora sejati}) saka awake dhewe; ora ana rong
objek kang beda kang saben-saben dadi perangan
saka liyane; perangan saka perangan sawijining
objek uga perangan saka objek kasebut. Gatekna
yen ing pangertèn iki "dadi perangan saka" mirip
kanthi "dadi subhimpunan saka", nanging ora mirip
kanthi "dadi anggota saka" kang ora refleksif
lan ora transitif.
\begin{align*}
& \lforall[x][\Atom{\Obj P}{x,x}] \\
& \lforall[x][\lforall[y][((\Part{x}{y} \land \Part{y}{x})
      \lif \eq[x][y])]] \\
& \lforall[x][\lforall[y][\lforall[z][((\Part{x}{y} \land
        \Part{y}{z}) \lif \Part{x}{z})]]]\\
\intertext{Kajaba iku, saben rong objek nduweni jumlah mereologis (objek kang
  nduweni loro-lorone minangka perangan, lan paling cilik ing bab iki).}  &
\lforall[x][\lforall[y][\lexists[z][\lforall[u][(\Part{z}{u} \liff
        (\Part{x}{u} \land \Part{y}{u}))]]]]
\end{align*}
Iki mung sawetara prinsip dhasar ngenani relasi
perangan kang digatekake dening para metafisikawan.
Nanging prinsip liyane enggal dadi angel dirumusake
utawa ditulis yen durung ngenalake sawetara relasi
kang ditemtokake luwih dhisik. Upamane, akeh
metafisikawan kang nggatekake mereologi uga
nganggep prinsip iki sah: kapan wae objek~$x$
nduweni perangan sejati~$y$, objek iku uga nduweni
perangan~$z$ kang ora nduweni perangan bareng
karo~$y$, lan gabungan $y$ lan $z$ iku~$x$.
\end{ex}

\end{document}
